\section{Generalized eigenvectors and generalized eigenspaces}

\begin{definition}
    Let $V$ be a finite-dimensional vector space over a field $F$ where $F$ is algebraically closed, and let
    \[
        T\in \operatorname{End}_F(V).
    \]
    Suppose that the characteristic polynomial of $T$ splits as
    \[
        f(x)=(x-\lambda_1)^{a_1}\cdots (x-\lambda_k)^{a_k},
    \]
    where $\lambda_1,\dots,\lambda_k$ are distinct. For each $i$, define
    \[
        V_{\lambda_i}:=\ker (T-\lambda_i I)^{a_i}.
    \]
    This is called the generalized eigenspace of $T$ corresponding to $\lambda_i$.
\end{definition}

\begin{theorem}
    With the notation above,
    \[
        V=V_{\lambda_1}\oplus \cdots \oplus V_{\lambda_k}.
    \]
\end{theorem}

\begin{proof}
    For each $i$, let
    \[
        P_i(x)=\frac{f(x)}{(x-\lambda_i)^{a_i}}
        =\prod_{j\ne i}(x-\lambda_j)^{a_j}.
    \]
    Since the polynomials $P_1(x),\dots,P_k(x)$ have greatest common divisor $1$, there exist polynomials
    \[
        Q_1(x),\dots,Q_k(x)\in F[x]
    \]
    such that
    \[
        Q_1(x)P_1(x)+\cdots+Q_k(x)P_k(x)=1.
    \]
    Replacing $x$ by $T$, we obtain
    \[
        Q_1(T)P_1(T)+\cdots+Q_k(T)P_k(T)=I_V.
    \]
    Hence every $v\in V$ can be written as
    \[
        v=\sum_{i=1}^k Q_i(T)P_i(T)v.
    \]
    Now
    \[
        (T-\lambda_i I)^{a_i}Q_i(T)P_i(T)
        =Q_i(T)f(T)=0,
    \]
    because $f(T)=0$ by the Cayley--Hamilton theorem. Therefore
    \[
        Q_i(T)P_i(T)v\in \ker (T-\lambda_i I)^{a_i}=V_{\lambda_i},
    \]
    and so
    \[
        V=V_{\lambda_1}+\cdots+V_{\lambda_k}.
    \]

    To show that the sum is direct, let
    \[
        v_1+\cdots+v_k=0,
        \qquad v_i\in V_{\lambda_i}.
    \]
    Fix $i$. Since $P_i(\lambda_i)\ne 0$, the polynomial $P_i(x)$ is relatively prime to $(x-\lambda_i)^{a_i}$. Hence there exist polynomials $A_i(x),B_i(x)\in F[x]$ such that
    \[
        A_i(x)P_i(x)+B_i(x)(x-\lambda_i)^{a_i}=1.
    \]
    Replacing $x$ by $T$ and applying to $v_i\in V_{\lambda_i}$, we get
    \[
        A_i(T)P_i(T)v_i=v_i,
    \]
    because $(T-\lambda_i I)^{a_i}v_i=0$. Thus $P_i(T)$ is invertible on $V_{\lambda_i}$.

    On the other hand, if $j\ne i$, then
    \[
        P_i(T)v_j=0,
    \]
    because $P_i(x)$ contains the factor $(x-\lambda_j)^{a_j}$ and
    \[
        (T-\lambda_j I)^{a_j}v_j=0.
    \]
    Applying $P_i(T)$ to the relation
    \[
        v_1+\cdots+v_k=0,
    \]
    we obtain
    \[
        P_i(T)v_i=0.
    \]
    Since $P_i(T)$ is invertible on $V_{\lambda_i}$, it follows that $v_i=0$. As $i$ was arbitrary, all $v_i$ are zero. Therefore the sum is direct.
\end{proof}

